\originalsection{18.112第一次历史考试（2006）}
\sourceinfo{Sigurdur Helgason课程公开材料, MIT OCW, CC BY-NC-SA 4.0. mid1prob.pdf的2008归档封面后, PDF p.2实印“Oct. 20, 2006”, 开卷; 官方解答mid1.pdf, pp.2--4. 以下保留卷面日期、5题及分值, 不改称2008试卷.}
\begin{exercise}\label{ass:m1-1}
\textbf{原题 M1.1（10分）。} 求方程 $z^3=-8\ii$ 的所有解.
\end{exercise}
\begin{solution}据官方解答翻译. $-8\ii=8\ee^{-\pi\ii/2}$, 所有根为 $2\ee^{\ii(-\pi/6+2k\pi/3)}$, $k=0,1,2$, 即 $\sqrt3-\ii,2\ii,-\sqrt3-\ii$. 模和辐角条件同时确定这三根, 因而没有遗漏.\end{solution}
\begin{exercise}\label{ass:m1-2}
\textbf{原题 M1.2（15分）。} 求正向 $\int_{|z-1|=1/2}(1-z)^{-3}\dd z$.
\end{exercise}
\begin{solution}据官方解答翻译. $(1-z)^{-3}=-(z-1)^{-3}$ 没有 $(z-1)^{-1}$ 项, 1处留数0, 积分为0. 用 Cauchy 二阶导数公式对常数函数1积分也给0.\end{solution}
\begin{exercise}\label{ass:m1-3}
\textbf{原题 M1.3（20分）。} 求 $\int_\gamma(\ee^z+z)/(z-2)\dd z$: 1) $\gamma$ 为 $|z|=1$; 2) $\gamma$ 为 $|z|=3$, 均取正向.
\end{exercise}
\begin{solution}据官方解答翻译. 1) 唯一奇点2在外, 内部解析, 积分0. 2) 2在内, Cauchy 公式给 $2\pi\ii(\ee^2+2)$.\end{solution}
\begin{exercise}\label{ass:m1-4}
\textbf{原题 M1.4（30分）。} $f$ 在全平面解析, 无穷远为非本性奇点, 证明 $f$ 是多项式.
\end{exercise}
\begin{solution}
据官方解答翻译并补全余项的全平面定义. 令 $g(w)=f(1/w)$ ($w\ne0$). 若0可去, $f$ 在无穷远有有限极限, 因而全平面有界, Liouville 给常数. 若0为 $h$ 阶极点, 在局部写 $g(w)=\sum_{k=1}^h B_kw^{-k}+\phi(w)$, $B_h\ne0$. 从 $g$ 减去有限主部后, $\phi$ 在 $\C\setminus\{0\}$ 解析且0可去, 故整. 当 $w\to\infty$, $g(w)\to f(0)$ 且主部趋0, 所以 $\phi$ 在无穷远有有限极限, 由 Liouville 是常数 $B_0$. 得 $f(z)=\sum_{k=0}^hB_kz^k$. 与第6章第3道自设练习的主结论重复, 该练习还要求多项式次数与无穷远极点阶数的对应; 本处保留官方题号及完整解答.
\end{solution}
\begin{exercise}\label{ass:m1-5}
\textbf{原题 M1.5（25分）。} $f$ 在全平面解析, 对某个 $n$ 及所有 $|z|>100$ 有 $|f(z)|<|z|^n$, 证明 $f$ 为多项式. 提示考察大的 $m$ 对应 $f^{(m)}(0)$.
\end{exercise}
\begin{solution}
据官方解答的 Cauchy 估计路线, 其模估计错误在此明确修正. $R>100$ 时 $M_R\le R^n$, 故 $|f^{(m)}(0)|\le m!M_R/R^m\le m!R^{n-m}$. 对任意非负整数 $m>n$ 令 $R\to\infty$, 得 $f^{(m)}(0)=0$. 整函数 Taylor 级数因此只剩有限项. 若通常取 $n\ge0$ 的整数, 次数至多 $n$; 原卷未另写 $n$ 类型, 实数 $n\ge0$ 时同样给次数至多 $\lfloor n\rfloor$, 若 $n<0$ 则连 $m=0$ 也为0, 即 $f\equiv0$.

官方PDF p.3在估计绝对值时写成复幂的围道积分, 并得到含 $1/(n-m)$ 的负上界, 不成立. 正确做法是积分 $|f(\zeta)|/|\zeta|^{m+1}$ 对弧长, 用圆周长度 $2\pi R$ 得上述非负上界.
\end{solution}
